\section{Justificación de la propuesta}

\noindent\textbf{Interés académico, científico e industrial.} La propuesta atiende una necesidad verificable del Laboratorio de Mecatrónica de EAFIT y, por extensión, de talleres de prototipado y microempresas de publicidad, para las cuales las alternativas comerciales son manuales y no repetibles, o bien sobredimensionadas y costosas \cite{shannon_line,mic_auto}. Automatizar el proceso lo vuelve reproducible y trazable, reduce el desperdicio y libera tiempo calificado. Científicamente, exige cuantificar la relación entre la historia térmica local y la recuperación elástica del PMMA, problema de física de polímeros con tratamiento constitutivo no trivial \cite{gilormini2010} e impacto directo sobre la precisión alcanzable.

\noindent\textbf{Áreas de conocimiento involucradas.} \emph{Transferencia de calor:} conducción transitoria a través del espesor, convección y radiación, número de Fourier y tiempo de penetración \cite{incropera2011}. \emph{Física de polímeros:} transición vítrea, viscoelasticidad no lineal, relajación de esfuerzos, \emph{crazing} y esfuerzos residuales \cite{osswald2012,gilormini2010,yan2023}. \emph{Mecánica de medios continuos:} flexión de láminas delgadas y recuperación elástica en estado gomoso. \emph{Instrumentación y metrología:} termometría por termopar e infrarroja, incertidumbre, calibración y metrología angular \cite{iso17025}, más la evaluación óptica de transparencia y defectos superficiales. \emph{Control, electrónica y sistemas embebidos:} identificación de modelo de primer orden con retardo, sintonización PID \cite{astrom2006}, conmutación de potencia, motores paso a paso e interfaz hombre-máquina.
